\documentclass{article}
\usepackage[a4paper,portrait,hmargin=10.5mm,top=8mm,bottom=8mm]{geometry}
\usepackage[T1]{fontenc}
\usepackage{lmodern}
\usepackage{amsmath,amssymb}
% Use scalable mathematical operators at the reference text sizes.
\input{omxlmex.fd}
\DeclareFontShape{OMX}{lmex}{m}{n}{<->lmex10}{}
\usepackage{xcolor,colortbl}
\usepackage{tabularx,array}
\pagestyle{empty}
\definecolor{navy}{HTML}{172C3C}
\definecolor{coral}{HTML}{BF604D}
\definecolor{rulecolor}{HTML}{A5AFB6}
\color{navy}
\setlength{\parindent}{0pt}
\setlength{\parskip}{0pt}
\setlength{\arrayrulewidth}{0.3pt}
% Central typography controls, explicitly applied in table cells.
\newcommand{\TitleFont}{\fontsize{25}{27}\selectfont\bfseries}
\newcommand{\KickerFont}{\fontsize{10}{12}\selectfont\bfseries}
\newcommand{\HeadingFont}{\fontsize{12.5}{15}\selectfont\bfseries}
\newcommand{\DefinitionFont}{\fontsize{9}{10.3}\selectfont}
\newcommand{\SymbolFont}{\fontsize{10}{11}\selectfont}
\newcommand{\ExampleFont}{\fontsize{8.8}{10.2}\selectfont}
\newcommand{\TableHeadFont}{\fontsize{8.5}{10}\selectfont\bfseries}
\newcommand{\FooterFont}{\fontsize{8.5}{10}\selectfont}
\DeclareMathSizes{10}{10}{7}{5}
\DeclareMathSizes{9}{9}{6.5}{5}
\DeclareMathSizes{8.8}{8.8}{6.3}{5}

\newlength{\SymbolWidth}
\newlength{\MeaningWidth}
\newlength{\RowPadding}
\setlength{\RowPadding}{0.7mm}
% Measure the complete cell (including wrapped text and tall mathematics),
% then reserve the same clearance above and below its actual bounds.
\newsavebox{\CellContent}
\newcommand{\paddedcell}[1]{%
  \sbox{\CellContent}{\begin{minipage}[t]{\linewidth}\raggedright #1\end{minipage}}%
  \raisebox{0pt}[\dimexpr\ht\CellContent+\RowPadding\relax]%
    [\dimexpr\dp\CellContent+\RowPadding\relax]{\usebox{\CellContent}}%
}
% Center only complete symbol rows, using the actual bounds of tall mathematics.
% Headers and the letter reference retain their original padded-cell layout.
\newcommand{\centeredcell}[1]{%
  \sbox{\CellContent}{\begin{minipage}[t]{\linewidth}\raggedright #1\end{minipage}}%
  \raisebox{\dimexpr(\dp\CellContent-\ht\CellContent)/2\relax}%
    [\dimexpr(\ht\CellContent+\dp\CellContent)/2+\RowPadding\relax]%
    [\dimexpr(\ht\CellContent+\dp\CellContent)/2+\RowPadding\relax]%
    {\usebox{\CellContent}}%
}
% Text-style arrays align relations without enlarging fractions or operators.
% Each r/E pair holds a left-hand side and its relation/right-hand side.
% E supplies an empty math atom to retain normal spacing before a relation.
\newcommand{\ExampleLineStretch}{1.12}
\newsavebox{\ExampleContent}
\newcommand{\exmath}[2][r@{}E]{\begingroup
  \renewcommand{\arraystretch}{\ExampleLineStretch}%
  \ifcsname NC@find@E\endcsname\else\newcolumntype{E}{>{{}}l}\fi%
  \sbox{\ExampleContent}{$\begin{array}[t]{@{}#1@{}}#2\end{array}$}%
  % Only overwide examples receive a small math-spacing correction.
  % This preserves the font size, baselines, and all fitting examples.
  \ifdim\wd\ExampleContent>\linewidth
    \thickmuskip=4mu\relax\medmuskip=3mu\relax
    \sbox{\ExampleContent}{$\begin{array}[t]{@{}#1@{}}#2\end{array}$}%
  \fi
  \usebox{\ExampleContent}\endgroup}
\newcommand{\symrow}[3]{%
  \centeredcell{\centering\SymbolFont\ensuremath{\displaystyle #1}}&%
  \centeredcell{\DefinitionFont #2}&\centeredcell{\ExampleFont #3}\\\hline
}
% Measure each example line so paired symbols and meanings share its baselines,
% including fractions and wrapped meanings. Spacing settings stay local.
\newsavebox{\PairFirst}
\newsavebox{\PairSecond}
\newsavebox{\PairThird}
% Override locally for compact measured rows; other sections retain their spacing.
\newcommand{\MeasuredLineStretch}{1}
\newcommand{\measuresymbol}[1]{}
\newcommand{\pairstrut}[1]{\rule[-\dp#1]{0pt}{\dimexpr\ht#1+\dp#1\relax}}
\newcommand{\pairlines}[3]{%
  \renewcommand{\arraystretch}{\MeasuredLineStretch}%
  \begin{tabular}[t]{@{}#1@{}}%
  \pairstrut{\PairFirst}#2\\%
  \pairstrut{\PairSecond}#3%
  \end{tabular}%
}
% Three subrows use the same measured struts and cell padding as paired rows.
\newcommand{\triplelines}[4]{%
  \renewcommand{\arraystretch}{\MeasuredLineStretch}%
  \begin{tabular}[t]{@{}#1@{}}%
  \pairstrut{\PairFirst}#2\\%
  \pairstrut{\PairSecond}#3\\%
  \pairstrut{\PairThird}#4%
  \end{tabular}%
}
\newcommand{\triplerow}[9]{%
  \noalign{\begingroup
  \renewcommand{\ExampleLineStretch}{\MeasuredLineStretch}%
  \global\setbox\PairFirst=\hbox{\ExampleFont\exmath{#7}\measuresymbol{#1}%
    \DefinitionFont\parbox[t]{\MeaningWidth}{\raggedright #4}}%
  \global\setbox\PairSecond=\hbox{\ExampleFont\exmath{#8}\measuresymbol{#2}%
    \DefinitionFont\parbox[t]{\MeaningWidth}{\raggedright #5}}%
  \global\setbox\PairThird=\hbox{\ExampleFont\exmath{#9}\measuresymbol{#3}%
    \DefinitionFont\parbox[t]{\MeaningWidth}{\raggedright #6}}%
  \endgroup}%
  \paddedcell{\centering\SymbolFont\triplelines{c}%
    {\ensuremath{\displaystyle #1}}{\ensuremath{\displaystyle #2}}{\ensuremath{\displaystyle #3}}}&%
  \paddedcell{\DefinitionFont\triplelines{l}%
    {\parbox[t]{\linewidth}{\raggedright #4}}%
    {\parbox[t]{\linewidth}{\raggedright #5}}%
    {\parbox[t]{\linewidth}{\raggedright #6}}}&%
  \paddedcell{\ExampleFont\renewcommand{\ExampleLineStretch}{\MeasuredLineStretch}%
    \exmath{\pairstrut{\PairFirst}#7\\\pairstrut{\PairSecond}#8\\\pairstrut{\PairThird}#9}}\\\hline
}
% Paired symbols, meanings, and examples; wide blocks can supply their columns.
\newcommand{\pairedrow}[6]{%
  \pairedrowwithcols{r@{}E@{\quad}r@{}E}{#1}{#2}{#3}{#4}{#5}{#6}}
\newcommand{\pairedrowwithcols}[7]{%
  \noalign{\begingroup
  \renewcommand{\ExampleLineStretch}{\MeasuredLineStretch}%
  \global\setbox\PairFirst=\hbox{\ExampleFont\exmath[#1]{#6}\measuresymbol{#2}%
    \DefinitionFont\parbox[t]{\MeaningWidth}{\raggedright #4}}%
  \global\setbox\PairSecond=\hbox{\ExampleFont\exmath[#1]{#7}\measuresymbol{#3}%
    \DefinitionFont\parbox[t]{\MeaningWidth}{\raggedright #5}}%
  \endgroup}%
  \paddedcell{\centering\SymbolFont\pairlines{c}{\ensuremath{\displaystyle #2}}{\ensuremath{\displaystyle #3}}}&%
  \paddedcell{\DefinitionFont\pairlines{l}%
    {\parbox[t]{\linewidth}{\raggedright #4}}%
    {\parbox[t]{\linewidth}{\raggedright #5}}}&%
  \paddedcell{\ExampleFont\renewcommand{\ExampleLineStretch}{\MeasuredLineStretch}%
    \exmath[#1]{%
    \pairstrut{\PairFirst}#6\\\pairstrut{\PairSecond}#7}}\\\hline
}
% Compact contexts use the existing palette and the table's example type.
\newcommand{\ContextHeading}{Example contexts}
\newcommand{\examplecontexts}[1]{%
  \begingroup
  \setlength{\fboxsep}{0.5mm}%
  \setlength{\fboxrule}{0.35pt}%
  \noindent\fcolorbox{coral}{coral!10}{%
    \begin{minipage}{\dimexpr\linewidth-2\fboxsep-2\fboxrule\relax}
      \ExampleFont\color{navy}%
      \textbf{\ContextHeading}\par
      #1\par
    \end{minipage}}\par
  \endgroup
  \vspace{0.7mm}%
}
% Optional contexts appear after the heading rule and before the table.
\newcommand{\SymbolColumnFraction}{.24}
\newcommand{\MeaningColumnFraction}{.32}
\newcommand{\sectionbox}[3][]{%
  {\HeadingFont\color{coral}#2\par}%
  \vspace{0.4mm}{\color{navy}\hrule height 0.5pt}\vspace{0.4mm}%
  \if\relax\detokenize{#1}\relax\else\examplecontexts{#1}\fi
  \begingroup
  \if\relax\detokenize{#1}\relax\else
    \renewcommand{\ExampleLineStretch}{1}%
  \fi
  \setlength{\SymbolWidth}{\SymbolColumnFraction\linewidth}%
  \setlength{\MeaningWidth}{\MeaningColumnFraction\linewidth}%
  \setlength{\tabcolsep}{0.7mm}%
  \renewcommand{\arraystretch}{1}%
  \arrayrulecolor{rulecolor}%
  \begin{tabularx}{\linewidth}{@{}>{\centering\arraybackslash\SymbolFont}p{\SymbolWidth}>{\raggedright\arraybackslash\DefinitionFont}p{\MeaningWidth}>{\raggedright\arraybackslash\ExampleFont}X@{}}
  \paddedcell{\centering\TableHeadFont Symbol}&\paddedcell{\TableHeadFont Meaning}&\paddedcell{\TableHeadFont Examples}\\\hline
  #3
  \end{tabularx}%
  \endgroup\par\vspace{0.8mm}%
}
% Dedicated two-column reference for letter names.
\newcommand{\letterrow}[2]{%
  \paddedcell{\centering\SymbolFont\ensuremath{#1}}&%
  \paddedcell{\DefinitionFont #2}\\\hline
}
\newcommand{\letterreference}[1]{%
  {\HeadingFont\color{coral}Letters and Greek symbols\par}%
  \vspace{0.4mm}{\color{navy}\hrule height 0.5pt}\vspace{0.4mm}%
  \begingroup
  \setlength{\tabcolsep}{0.7mm}%
  \renewcommand{\arraystretch}{1}%
  \arrayrulecolor{rulecolor}%
  \begin{tabularx}{\linewidth}{@{}>{\centering\arraybackslash\SymbolFont}p{.29\linewidth}>{\raggedright\arraybackslash\DefinitionFont}X@{}}
  \paddedcell{\centering\TableHeadFont Symbol}&\paddedcell{\TableHeadFont Name / description}\\\hline
  #1
  \end{tabularx}%
  \endgroup\par\vspace{0.8mm}%
}

% Half-width tables share baseline alignment, with local widths and padding.
% General notation retains its original complete-row centering and spacing.
\newcommand{\compacttables}{%
  \renewcommand{\measuresymbol}[1]{\SymbolFont$\displaystyle ##1$}%
  \setlength{\RowPadding}{0.18mm}%
  \renewcommand{\SymbolColumnFraction}{.22}%
  \renewcommand{\MeaningColumnFraction}{.36}%
  \renewcommand{\pairedrow}[6]{\pairedrowwithcols{r@{}E}{##1}{##2}{##3}{##4}{##5}{##6}}%
  \renewcommand{\symrow}[3]{%
    \paddedcell{\centering\SymbolFont\ensuremath{\displaystyle ##1}}&%
    \paddedcell{\DefinitionFont ##2}&\paddedcell{\ExampleFont ##3}\\\hline}%
}
% Letter names use the same measured subrows as the three-column tables.
\newcommand{\letterpair}[4]{%
  \noalign{\global\setbox\PairFirst=\hbox{\SymbolFont$#1$\DefinitionFont #3}%
    \global\setbox\PairSecond=\hbox{\SymbolFont$#2$\DefinitionFont #4}}%
  \paddedcell{\centering\SymbolFont\pairlines{c}{$#1$}{$#2$}}&%
  \paddedcell{\DefinitionFont\pairlines{l}{#3}{#4}}\\\hline
}
\newcommand{\lettertriple}[6]{%
  \noalign{\global\setbox\PairFirst=\hbox{\SymbolFont$#1$\DefinitionFont #4}%
    \global\setbox\PairSecond=\hbox{\SymbolFont$#2$\DefinitionFont #5}%
    \global\setbox\PairThird=\hbox{\SymbolFont$#3$\DefinitionFont #6}}%
  \paddedcell{\centering\SymbolFont\triplelines{c}{$#1$}{$#2$}{$#3$}}&%
  \paddedcell{\DefinitionFont\triplelines{l}{#4}{#5}{#6}}\\\hline
}

\begin{document}
\begingroup
{\TitleFont\color{coral}MATHEMATICS\hfill\KickerFont POSTER 1\par}
\vspace{1.3mm}{\color{navy}\hrule height 0.8pt}\vspace{1.8mm}
{\KickerFont GENERAL NOTATION, SETS AND PROBABILITY\par}\vspace{2mm}
\begingroup
% Local column and row spacing keeps the expanded reference on Poster 1.
\setlength{\RowPadding}{0.05mm}
\renewcommand{\MeasuredLineStretch}{0.75}
\renewcommand{\SymbolColumnFraction}{.23}
\renewcommand{\MeaningColumnFraction}{.41}
\sectionbox{General notation}{
\triplerow{=}{\neq}{\approx}{Equal}{Not equal}{Approximately}
  {2+3&=5\qquad \qquad 4-1=3}{2&\neq3\qquad \qquad 5+1\neq7}{\pi&\approx3.14\> \> \> \quad \qquad \sqrt2\approx1.41}

\pairedrow{<,\ >}{\nless,\ \ngtr}
{\qquad Less than \qquad \qquad Greater than}
{Not Less than \qquad Not Greater than}
{2&<3 \qquad &5&>4}{2&\nless2 \qquad &4&\ngtr5}

\pairedrow{\leq,\ \geq}{\nleq,\ \ngeq}
{\qquad Less/Greater than or equal to}
{ Not Less/Greater than or equal to}
{2&\leq2 \qquad &5&\geq4}{3&\nleq2 \qquad &4&\ngeq5}

\symrow{\pm}{Plus or minus: two choices}{\exmath{3\pm1&=4\text{ or }2\qquad -2\pm3=1\text{ or }-5}}
\symrow{ab=a(b)=a\cdot b=a\times b}
{Multiplication: \\ juxtaposition, parentheses, dot, or cross.}
{\exmath{3x&=3\cdot x=3\times x\\
		2(4)&=2\cdot4=2\times4}}

\symrow{\frac{a}{b}=a/b=a\div b}
{Division: \\ fraction, slash, or division sign ($b\neq0$).}
{\exmath[r@{}E@{}E@{}E]{\frac{12}{3}&=12/3&=12\div3&=4\\
		\frac{x}{5}&=x/5&=x\div 5&=0.2x}}
\symrow{|a|}{Absolute value: distance from $a$ to $0$}{\exmath{|-5|&=5 \qquad\quad  |3|=3\qquad\quad  |0|=0}}
\pairedrowwithcols{r@{}E@{\;}r@{}E}{a^2}{a^3}
  {Squared; multiply $a$ by itself.}
  {Cubed; multiply three factors of $a$.}
  {3^2&=3\cdot3=9 &(-2)^2&= (-2)(-2) = 4}
  {2^3&=2\cdot2\cdot2=8 \quad &(-2)^3&=(-2)(-2)(-2)=-8}
\pairedrowwithcols{r@{}E}{\sqrt{a}=a^{1/2}}{\sqrt[3]{a}=a^{1/3}}
  {Square root \quad \>  $a\geq 0$}
  {Cube root \qquad \! $a \in \mathbb{R}$}
  {\sqrt9&=3\text{ (not }\pm3\text{)}}
  {\sqrt[3]{-8}&=-2}
\pairedrowwithcols{r@{}E@{\;}r@{}E}
	{a^n}{\sqrt[n]{a}=a^{1/n}}
	{Exponentiation: $n$ factors of $a$ }
	{$n$th root.}
	{2^4&=2\cdot2\cdot2\cdot2=16 \quad &(-2)^5&=-32}
	{\sqrt[4]{16}& = 16^{\frac{1}{4}}=2 \qquad &\sqrt[5]{-32}&=-2 = (-32)^{\frac{1}{5}}}
}
\endgroup
\vspace{0mm}
\noindent\begin{minipage}[t]{.49\linewidth}
\compacttables
\renewcommand{\MeaningColumnFraction}{.34}
\sectionbox{Sets and number systems}{
\symrow{\{a,b,c\}}{Set: ignores order and repetition.}{\exmath{\{2,4\}&=\{4,2\}\\\{2,2,4\}&=\{2,4\}}}
\pairedrow{\in}{\notin}{Element of}{Not an element of}{2&\in\{1,2\}}{3&\notin\{1,2\}}
\symrow{\{x:p(x)\}}{$x\in A$ such that $p(x)$ holds.}{\exmath[l]{\{x\in\mathbb R:x>0\}\\\text{Positive real numbers}}}
\pairedrow{\varnothing}{|A|}{Empty set: no elements.}{Number of elements in finite $A$.}{|\varnothing|&=0}{|\{2,4,6\}|&=3}
\pairedrow{A\subseteq B}{A\subsetneq B}{Subset: all of $A$ in $B$; may equal $B$.}{Proper subset: $A\subseteq B$, $A\ne B$.}{\{1,2\}&\subseteq\{1,2\}}{\{1\}&\subsetneq\{1,2\}}
\pairedrow{A\cup B}{A\cap B}{Union: in either or both.}{Intersection: in both.}{\{1\}\cup\{2\}&=\{1,2\}}{\{1,2\}\cap\{2\}&=\{2\}}
\symrow{A\setminus B}{Difference: in $A$, outside $B$.}{\exmath{\{1,2\}\setminus\{2\}&=\{1\}}}
\pairedrow{\mathbb N}{\mathbb Z}{Naturals}{Integers}{\mathbb N&=\{1,2,3,\ldots\}}{Z&=\{\ldots, -2, -1, 0, 1,2,\ldots\}}
\pairedrow{\mathbb Q}{\mathbb R}{Rationals: $p/q$; $p,q\in\mathbb Z$, $q\ne0$.}{Reals: the number line.}{\tfrac12,-3&\in\mathbb Q}{\sqrt2,\pi,\tfrac12,-3&\in\mathbb R}
\symrow{\mathbb R\setminus\mathbb Q}{Irrationals}{\exmath{\sqrt2,\pi&\in\mathbb R\setminus\mathbb Q}}
\pairedrowwithcols{l@{}E}{(a,b)}{[a,b]}{Open: excludes endpoints.}{Closed: includes endpoints.}{x\in(1,3):&1<x<3}{x\in[1,3]:&1\leq x\leq3}
\pairedrowwithcols{l@{}E}{[a,b)}{(a,b]}{Includes left endpoint only.}{Includes right endpoint only.}{2\in[2,5),&5\notin[2,5)}{2\notin(2,5],&5\in(2,5]}
\pairedrowwithcols{l}{\infty}{-\infty}{Unbounded above.}{Unbounded below.}{[0,\infty)}{(-\infty,0]}
}
{\FooterFont Intervals assume $a<b$. Infinity is not a real number; use parentheses there. Some texts include $0$ in $\mathbb N$; check the convention. Dots: see Poster 2. Bars: absolute value for numbers, cardinality for finite sets.\par}
\end{minipage}\hfill
\begin{minipage}[t]{.485\linewidth}
\compacttables
\sectionbox[{\textbf{Fair die:} $\Omega=\{1,2,3,4,5,6\}$.\par
$A=\{1,3,5\}$ (odd), $B=\{1,2\}$, $C=\{2,4,6\}$.\par
Die examples below use these events.\par
\textbf{Independent fair coin tosses:} $E$: first is heads; $F$: second is heads.\par}]{Probability and events}{
\pairedrow{\Omega}{\omega}{Sample space: all possible outcomes.}{Outcome: one possible result.}{\Omega&=\{1,2,3,4,5,6\}}{\omega&=3\in\Omega}
\symrow{A\subseteq\Omega}{Event: a set of outcomes.}{\exmath{A&=\{1,3,5\}}}
\symrow{P(A)}{Probability of $A$: between $0$ and $1$.}{\exmath{0&\leq P(A)\leq1\\P(\Omega)&=1\\P(\varnothing)&=0}}
\symrow{\dfrac{|A|}{|\Omega|}}{Equals $P(A)$ for finite, equally likely outcomes only.}{\exmath{P(A)&=\tfrac36=\tfrac12}}
\symrow{A^c}{Complement: outcomes in $\Omega$ outside $A$.}{\exmath{A^c&=\{2,4,6\}\\P(A^c)&=1-P(A)=\tfrac12}}
\pairedrow{A\cup B}{A\cap B}{Event $A$ or $B$, or both.}{Events $A$ and $B$ together.}{A\cup B&=\{1,2,3,5\}}{A\cap B&=\{1\}}
\symrow{A\cap C=\varnothing}{Mutually exclusive: cannot occur together.}{\exmath{A\cap C&=\varnothing}}
\symrow{P(A\mid B)}{Conditional probability: $A$ given $B$; $P(B)>0$.}{\exmath{P(A\mid B)&=\tfrac{P(A\cap B)}{P(B)}\\&=\tfrac{1/6}{2/6}=\tfrac12}}
\symrow{P(E\cap F)}{Independent events: multiply their probabilities.}{\exmath{P(E\cap F)&=P(E)P(F)\\&=\tfrac12\cdot\tfrac12=\tfrac14}}
\symrow{P(A\mid B)}{Independence iff this equals $P(A)$; $P(B)>0$.}{\exmath{P(E\mid F)&=P(E)=\tfrac12}}
}
\vspace{1.5mm}
{\DefinitionFont\textbf{Keep the ideas separate.}\par
Mutually exclusive events of positive probability are not independent: one occurring rules out the other.\par
\vspace{2mm}
The bar in $P(A\mid B)$ means ``given''; in set-builder notation it means ``such that''. Context determines the reading.\par}
\end{minipage}\par
\vfill
{\color{navy}\hrule height 0.5pt}\vspace{1.5mm}
{\FooterFont Symbols acquire meaning from their mathematical context.\par}
\endgroup
\newpage
\begingroup
{\TitleFont\color{coral}MATHEMATICS\hfill\KickerFont POSTER 2\par}
\vspace{1.3mm}{\color{navy}\hrule height 0.8pt}\vspace{1.8mm}
{\KickerFont FUNCTIONS, VARIABLES, SEQUENCES AND SUMS\par}\vspace{3mm}
\noindent\begin{minipage}[t]{.49\linewidth}
\compacttables
\sectionbox[{\textbf{Example:} $f:\mathbb R\to\mathbb R$, $f(x)=x^2$.\par
Set and interval notation: see Poster 1.\par}]{Function notation}{
\symrow{f}{Function: one output per domain input.}{\exmath{f(x)&=x^2}}
\symrow{f(x)}{Value at $x$; parentheses mark input, not multiplication.}{\exmath{f(3)&=3^2=9}}
\symrow{f:A\to B}{Maps domain $A$ to codomain $B$.}{\exmath[l]{f:\mathbb R\to\mathbb R}}
\pairedrow{A}{B}{Domain: allowed inputs.}{Codomain: declared output target.}{A&=\mathbb R}{B&=\mathbb R}
\symrow{x\mapsto f(x)}{Maps values; $\to$ connects sets.}{\exmath[r@{}E]{x&\mapsto x^2\\3&\mapsto9}}
\symrow{\operatorname{dom}(f)}{Domain: allowed inputs.}{\exmath{\operatorname{dom}(f)&=\mathbb R}}
\symrow{\operatorname{range}(f)}{Range: attained outputs; within codomain.}{\exmath{\operatorname{range}(f)&=[0,\infty)}\par No negative outputs.}
}
\vspace{1.5mm}
\letterreference{
\letterrow{a,b,c;\ x,y,z}{Common lowercase Latin letters}
\letterrow{A,B,C}{Common uppercase Latin letters}
\lettertriple{\alpha}{\beta}{\gamma}{Alpha}{Beta}{Gamma}
\letterpair{\delta}{\Delta}{Delta}{Uppercase delta}
\letterpair{\theta}{\Theta}{Theta}{Uppercase theta}
\letterpair{\lambda}{\Lambda}{Lambda}{Uppercase lambda}
\lettertriple{\mu}{\sigma}{\Sigma}{Mu}{Sigma}{Uppercase sigma}
\letterpair{\omega}{\Omega}{Omega}{Uppercase omega (Poster 1)}
\letterrow{\pi}{Pi; here, the circle constant $\pi\approx3.14$.}
}
{\DefinitionFont Names do not fix meanings. $\Sigma$ is a letter; $\sum$ is the related summation operator (right), not an ordinary variable.\par}
\vspace{3mm}
\renewcommand{\ContextHeading}{Letter roles}
\sectionbox[{Roles depend on context: uppercase often names sets; Greek often names angles or numbers. Conventions are not permanent types.\par}]{Variables}{
\pairedrowwithcols{l}{x}{m,b}{Variable: may vary in context.}{Constants: fixed in context.}{x\in\{-2,0,3.5\}}{y=mx+b}
\pairedrowwithcols{l}{x}{y}{Input for a line.}{Output for that line.}{y=mx+b}{m,b\text{ fixed}}
\triplerow{n}{\theta}{A}{May count objects.}{May name an angle in degrees (${}^\circ$).}{May name a set.}{n&\in\mathbb N}{\theta&=30^\circ}{A&=\{1,2\}}
\symrow{\lambda}{May name a number.}{\exmath{\lambda&=2}}
}
{\FooterFont A letter may be variable or constant. Values $-2$, $0$, $3.5$ are separate choices for $x$, not simultaneous equalities.\par}
\end{minipage}\hfill
\begin{minipage}[t]{.485\linewidth}
\compacttables
\renewcommand{\SymbolColumnFraction}{.25}
\renewcommand{\MeaningColumnFraction}{.34}
\renewcommand{\ContextHeading}{Reading guide}
\sectionbox[{A \textbf{sequence} is an ordered list of terms; a \textbf{series} is a sum of sequence terms. Here, parentheses show sequences; braces show sets.\par
\textbf{Sequence examples:} $a_i=2i$, starting at $i=1$.\par}]{Lists, sequences and subscripts}{
\pairedrowwithcols{l}{\ldots}{\cdots}{Omitted list entries.}{Omitted terms in a sum.}{1,2,3,\ldots}{a_1+a_2+\cdots+a_n}
\pairedrowwithcols{l}{1,2,\ldots,n}{1,2,3,\ldots}{Finite list: ends at $n$.}{Infinite list: no final term.}{1,2,3,\ldots,5\ (n=5)}{1,2,3,4,5,\ldots}
\symrow{a_i}{Term at position or label $i$.}{\exmath{a_3&=2(3)=6}}
\pairedrow{a_1}{a_n}{Term at index $1$.}{Term at index $n$.}{a_1&=2}{a_n&=2n}
\pairedrow{(a_i)_{i=1}^{4}}{(a_i)_{i=1}^{\infty}}{Finite: four terms.}{Infinite: no last term.}{(a_i)_{i=1}^{4}&=(2,4,6,8)}{(a_i)_{i=1}^{\infty}&=(2,4,6,\ldots)}
\pairedrowwithcols{l}{(2,4)}{(2,2,4)}{Order matters.}{Repetition matters.}{(2,4)\ne(4,2)}{(2,2,4)\ne(2,4)}
\symrow{A_1,A_2}{Subscripts may label cases.}{Alternatives $A_1$ and $A_2$.}
}
{\FooterFont Dots follow an established pattern; alone they do not define a unique rule. The finite list $1,2,3,\ldots,n$ assumes integer $n\geq3$. An index is not multiplication or an exponent. Indices may start at $1$, $0$, or another integer: state the start. Sets ignore order and repetition (Poster 1).\par}
\vspace{3mm}
\renewcommand{\ContextHeading}{Index convention}
\renewcommand{\SymbolColumnFraction}{.23}
\renewcommand{\MeaningColumnFraction}{.31}
\sectionbox[{Indices run through consecutive integers between the bounds, \textbf{inclusive}. Here $n\in\mathbb N$, $n\geq1$, and $a_i=2i$ in the sequence examples.\par}]{Summation notation}{
\symrow{\displaystyle\sum_{i=1}^{n}a_i}{Sum of the indicated terms.}{\exmath[l]{a_1+a_2+\cdots+a_n}}
\pairedrowwithcols{l}{i}{a_i}{Index: local to sum.}{Summand: expression summed.}{i=1,2,\ldots,n}{a_i=2i}
\pairedrowwithcols{l}{1}{n}{Lower bound: first index.}{Upper bound: last index.}{\text{Start at }i=1.}{\text{End at }i=n.}
\pairedrow{a_n}{S_n}{Single term at $n$.}{Sum of first $n$ terms; start at $1$.}{a_4&=8}{S_n&=\sum_{i=1}^{n}a_i}
\symrow{S_4}{Four terms for $a_i=2i$.}{\exmath{S_4&=2+4+6+8=20}}
\pairedrow{\displaystyle\sum_{i=2}^{4}3}{\displaystyle\sum_{k=1}^{3}k^2}{Constant; bounds $2$ to $4$.}{Squares; bounds $1$ to $3$.}{\sum_{i=2}^{4}3&=3+3+3=9}{\sum_{k=1}^{3}k^2&=1+4+9=14}
\symrow{i,\ k}{Index names may change consistently.}{\exmath{\sum_{i=1}^{4}i&=\sum_{k=1}^{4}k\\&=1+2+3+4=10}}
}
{\FooterFont Rename the local index in both bounds and summand. Centered dots use the same omission convention as list dots.\par}
\end{minipage}\par
\vfill
{\color{navy}\hrule height 0.5pt}\vspace{1.5mm}
{\FooterFont Symbols acquire meaning from their mathematical context.\par}
\endgroup
\end{document}
